% Modern TATA §1.5 companion; NOT the unheld Thatcher--Wright 1968 presentation.
\section{Finite Algebras and Context Congruence for Ranked Trees}
\label{sec:p09-finite-tree-algebra}

The finite-algebra viewpoint asks what finite information about a subtree
suffices in every surrounding tree. The source for this treatment is the
fixed 18 November 2008 release of \emph{Tree Automata Techniques and
Applications}, especially its tree Myhill--Nerode theorem, algebraic
corollary, and minimization algorithm in \S1.5
\parencite[pp.~35--37]{ComonEtAl2007}.
It supplements the existing finite-tree and WSkS constructions in
\cref{sec:p09-presentation-n-c-015,thm:p09-finite-tree-automata};
it is not their replacement or a reconstruction attributed to the
unheld 1968 Thatcher--Wright article.

\subsection{The exact ranked-tree input and algebra}
\label{sec:p09-tree-algebra-input}

Let \(F=\bigsqcup_{n\geq0}F_n\) be a finite ranked alphabet, with
\(F_0\ne\varnothing\), as in TATA's preliminaries. A ground term in
\(T(F)\) is a finite ordered tree: a node labelled \(f\in F_n\)
has exactly the children \(1,\ldots,n\), and a constant is a leaf.
Order of the children belongs to the input, not to an arbitrary drawing
\parencite[pp.~15--17, 20--21]{ComonEtAl2007}.
The previous notation \(T_\Sigma\) agrees after setting \(\Sigma=F\);
TATA's final-state set \(Q_f\) must not be confused with this alphabet.

A one-hole ground context \(C\in C(F)\) is a term with exactly one
occurrence of a distinguished hole \(x\), all other leaves being constants
of \(F\). The trivial context is \(x\) itself. Write \(C[t]\) for
replacement of its hole by \(t\). Context composition \(C[D]\) inserts
the one-hole context \(D\) into \(C\) and again has exactly one hole.
A context's hole depth is the number of edges from its root to that hole.

\begin{hbdefinition}{def:p09-finite-f-algebra}{Finite ranked-signature algebra and recognition}
An \(F\)-algebra \(A\) has a nonempty carrier, denoted \(|A|\) in the
source's algebra notation, and a total operation
\(f^A:|A|^n\to|A|\) for every \(f\in F_n\).
The case \(n=0\) is a designated element \(f^A\in|A|\).
Here \(|A|\) denotes the carrier, not its cardinality; write
\(\operatorname{card}(|A|)\) for the latter.
The algebra is finite when its carrier is finite.
A homomorphism \(h:A\to B\) preserves every operation, including
each designated constant.

The term algebra has carrier \(T(F)\) and operations that form terms.
A finite-algebra recognition of \(L\subseteq T(F)\) consists of
a finite \(F\)-algebra \(A\), a homomorphism
\(\varphi:T(F)\to A\), and a designated subset \(A'\subseteq|A|\)
such that \(L=\varphi^{-1}(A')\).
The designated subset is a recognition condition, not an additional
operation in the signature.
\end{hbdefinition}

\begin{hblemma}{thm:p09-tree-algebra-evaluation}{proved=annex:p09.tree-algebra.evaluation}{Unique ground-term evaluation and accessible subalgebra}
Every \(F\)-algebra has a unique ground-term homomorphism
\(\varphi_A:T(F)\to A\), given recursively by
\[
 \varphi_A(f(t_1,\ldots,t_n))
       =f^A(\varphi_A(t_1),\ldots,\varphi_A(t_n)).
\]
Its image is the nonempty subalgebra generated by the constant
interpretations. For finite \(A\), that image is computed by starting
with those constants and repeatedly closing under all operations.
Every obtained element has a finite term witness.
\end{hblemma}

Totality makes every ground term evaluable. A finite algebra is therefore
a deterministic complete bottom-up machine: its carrier is the state
set, its operations are transition functions, and \(A'\) is the final
set. Conversely a deterministic complete machine gives just such an
algebra. A nondeterministic machine is first determinized by the existing
subset construction of \cref{sec:p09-finite-tree-determinization}.
That construction retains the empty subset; deleting it would lose
totality and invalidate this identification.

\subsection{Congruence and the power of surrounding contexts}
\label{sec:p09-tree-context-congruence}

\begin{hbdefinition}{def:p09-tree-context-congruence}{Source-typed congruence and contextual indistinguishability}
An equivalence relation \(\sim\) on \(T(F)\) is a congruence when
\(u_i\sim v_i\) for every \(1\leq i\leq n\) implies
\(f(u_1,\ldots,u_n)\sim f(v_1,\ldots,v_n)\).
Its index is the number of its equivalence classes.
It saturates \(L\) when every class is either contained in \(L\)
or disjoint from \(L\).

Retain the source notation \(\equiv_L\), with the explicit local type
of a tree-language congruence, not elementary equivalence:
\[
 u\equiv_L v
 \quad\Longleftrightarrow\quad
 \text{for every }C\in C(F),\
           (C[u]\in L\ \Longleftrightarrow\ C[v]\in L).
\]
The two displayed biconditionals are metalinguistic truth equivalences.
They do not quantify over arbitrary nonfinite trees or contexts with
multiple independent holes.
\end{hbdefinition}

Checking only \(u\in L\) versus \(v\in L\) is insufficient:
terms can have the same immediate acceptance but differ after a parent
operation. The trivial context ensures that contextual equivalence implies
the same immediate acceptance. Nontrivial contexts ensure that the
identification remains safe in every enclosing tree.

\begin{hblemma}{thm:p09-tree-congruence-contexts}{proved=annex:p09.tree-algebra.contexts}{Operator compatibility and one-hole contexts}
An equivalence relation on \(T(F)\) is a congruence exactly when it is
preserved by every one-hole context. The relation \(\equiv_L\) is a
congruence saturating \(L\). Every congruence that saturates \(L\)
refines \(\equiv_L\).
\end{hblemma}

The last assertion is the crucial direction: the syntactic congruence is
the coarsest safe identification, not an arbitrarily chosen equivalence
of finitely many samples. Its definition refers to all finite contexts
and need not give an effective finite computation for an arbitrary
language \(L\).

\begin{hbtheorem}{thm:p09-tree-algebra-recognition}{proved=annex:p09.tree-algebra.recognition}{Finite-algebra and finite-congruence characterizations}
For \(L\subseteq T(F)\), the following are equivalent:
\begin{enumerate}
\item \(L\) is recognized by a finite bottom-up tree automaton.
\item \(L\) is recognized by a finite \(F\)-algebra.
\item A finite-index congruence saturates \(L\).
\item The congruence \(\equiv_L\) has finite index.
\end{enumerate}
When these hold, the quotient algebra has carrier
\(Q_{\min}=T(F)/{\equiv_L}\), operations
\[
 f^{A_L}([t_1],\ldots,[t_n])=[f(t_1,\ldots,t_n)],
\]
and designated set \(Q_{\min,f}=\{[t]:t\in L\}\).
Its ground-term evaluation sends \(t\) to \([t]\), so it recognizes
exactly \(L\).
\end{hbtheorem}

This is the source's three-part Myhill--Nerode characterization together
with its finite-algebra corollary, with the homomorphism direction made
explicit \parencite[p.~36]{ComonEtAl2007}.
Well-defined quotient operations are a proved requirement:
choosing different representatives cannot change the resulting class.
Merely naming a quotient would not establish this.

\subsection{Minimality and the complete-machine boundary}
\label{sec:p09-tree-algebra-minimal}

\begin{hbtheorem}{thm:p09-tree-algebra-minimal}{proved=annex:p09.tree-algebra.minimal}{Canonical quotient and minimum complete deterministic recognition}
For any finite-algebra recognition \((A,\varphi,A')\) of \(L\), let
\(R=\varphi(T(F))\). There is a unique surjective homomorphism
\[
 \pi:R\to A_L,\qquad \pi(\varphi(t))=[t],
\]
whose inverse image of the quotient's designated set is \(A'\cap R\).
Thus every reachable complete deterministic recognizer has at least
\(\operatorname{card}(Q_{\min})\) states. Equality gives an isomorphism
preserving operations and designated sets. Consequently the minimum
complete deterministic recognizer is unique up to state renaming.
\end{hbtheorem}

Reachability matters. Unused states in a supplied algebra carry no
information about evaluated ground terms and need not survive.
The source's Corollary~1.5.1 and algorithm MIN use the canonical quotient
and a reduced machine, respectively \parencite[p.~37]{ComonEtAl2007}.
Here reduced means every state has a ground-term witness.

Completeness also matters. Over the nullary alphabet \(F=\{a,b\}\),
the language \(\{a\}\) has a one-state partial deterministic recognizer
with only \(a\to q\), \(q\) final. A complete recognizer must distinguish
\(a\) from \(b\), and therefore needs two states.
The algebraic minimum statement is explicitly about total operations;
it is not an unqualified minimum over partial machines.

TATA assumes at least one constant, whereas the earlier handbook tree
definition also permits alphabets with no ground terms. If that hypothesis
is dropped, there are no input terms; the source quotient's ordinary
nonempty-carrier/minimality presentation needs a separate empty-input
convention. The present source treatment retains \(F_0\ne\varnothing\)
rather than silently extending its uniqueness claim to that regime.

\subsection{Computing the congruence by refinement}
\label{sec:p09-tree-algebra-refinement}

Start with a deterministic complete machine and remove unreachable states
by the evaluation closure construction. Write the remaining machine as
\((Q,F,Q_f,\delta)\).
The source's MIN algorithm begins with acceptance classes
\(\{Q_f,Q\setminus Q_f\}\), with empty blocks discarded.
Let \(P_0\) be the corresponding equivalence relation.

Given \(P_j\), define \(qP_{j+1}r\) when \(qP_jr\) and, for every
symbol \(f\in F_n\), every position \(1\leq i\leq n\), and every
tuple of other states,
\[
 \delta_f(q_1,\ldots,q_{i-1},q,q_{i+1},\ldots,q_n)
 \ P_j\
 \delta_f(q_1,\ldots,q_{i-1},r,q_{i+1},\ldots,q_n).
\]
Nullary operations contribute no variable position to this test.
Use the old \(P_j\) throughout one round; mixing updates within a round
would no longer be this displayed recurrence.

Stop when no class splits. The quotient transitions apply \(\delta_f\)
to representatives and return their final class; a class is final
exactly when its members lie in \(Q_f\).

\begin{hbtheorem}{thm:p09-tree-algebra-refinement}{proved=annex:p09.tree-algebra.refinement}{Termination, exact context invariant, and minimal output}
Each \(P_j\) is an equivalence relation. States are \(P_j\)-equivalent
exactly when every one-hole context of hole depth at most \(j\)
gives them the same final/nonfinal result.
The process terminates after at most
\(\operatorname{card}(Q)-\operatorname{card}(Q/P_0)\) strict splitting
rounds. At stability, the relation is a congruence, its quotient is
well defined, and the output is the canonical minimum complete
deterministic recognizer of the same language.
\end{hbtheorem}

This fills out the correctness argument for MIN at the source's
finite-input scope. It is not an assertion that every language supplied
only as a predicate has an effective syntactic congruence, or that this
straightforward implementation is an optimal minimization algorithm.

\subsection{Non-word examples and concrete refinements}
\label{sec:p09-tree-algebra-examples}

Use \(F=\{a^{(0)},b^{(0)},f^{(2)}\}\).
These worked examples are handbook-owned uses of the source constructions,
not examples attributed to the unheld 1968 paper.
For parity of \(b\)-leaves, take carrier \(\{0,1\}\),
\(a^A=0\), \(b^A=1\), \(f^A(u,v)=u+v\pmod2\), and designated
set \(\{0\}\).
The evaluation of \(f(b,f(a,b))\) gives \(1\) at the left leaf,
\(0,1\) at the right subtree's leaves, \(1\) at that subtree,
and \(0\) at the root. It is an accepting finite branching run.
Both states are reachable by constants.

For a supplied three-state recognizer use carrier \(\{0,1,2\}\),
constant values \(0,1\), operation \(f(u,v)=u+v\pmod3\), and
designated set \(\{0\}\). It recognizes divisibility of the number
of \(b\)-leaves by three, not parity. Its initial partition is
\(\{0\},\{1,2\}\). Substitution into the context \(f(x,b)\)
separates \(1\) from \(2\): the outputs are \(2\) and \(0\).
MIN therefore splits the latter block. This is a worked refinement,
not an invalid claim that a three-state mod-three recognizer reduces
to the two-state parity machine.

To exhibit actual redundant states for parity, instead use
\(Q=\{0,1,2,3\}\), constants \(a\mapsto0,b\mapsto1\),
addition modulo four at \(f\), and \(Q_f=\{0,2\}\).
All four states are reachable. States of the same parity remain
indistinguishable under every operation, so MIN yields precisely
\(\{0,2\},\{1,3\}\).

Order is independently visible in the leftmost-leaf algebra:
\(a^B=0,b^B=1,f^B(u,v)=u\), designated set \(\{0\}\).
It accepts \(f(a,b)\) and rejects \(f(b,a)\).
Thus identical leaf counts and an unordered diagram do not determine
the result for every ranked-tree algebra.

\begin{hbtheorem}{thm:p09-tree-algebra-cases}{proved=annex:p09.tree-algebra.cases}{The displayed runs, refinements, minima, and partial-machine countercase}
The parity algebra evaluates a tree to its number of \(b\)-leaves modulo
two and is minimal among complete deterministic recognizers of its
language. The modulo-three refinement splits \(1,2\); the modulo-four
parity recognizer minimizes to two classes.
The leftmost algebra evaluates to the label type of the leftmost leaf.
The nullary-language partial/complete countercase requires respectively
one and two states.
\end{hbtheorem}

The previously proved copying-language obstruction in
\cref{sec:p09-presentation-n-c-015} has a contextual reading:
for distinct \(t,u\), the context \(f(t,x)\) accepts \(t\) in the
language of equal-child roots and rejects \(u\). It therefore has
infinitely many context classes. This is the same obstruction, not a
new source acquisition or a second credit for its existing proof.

\subsection{The logic and decision-transfer seam}
\label{sec:p09-tree-algebra-logic-seam}

\begin{hbcorollary}{thm:p09-tree-algebra-mso}{proved=annex:p09.tree-algebra.mso}{Finite algebra recognition and finite-domain MSO}
Under the fixed finite ranked alphabet and the label/child signature of
\cref{thm:p09-finite-tree-automata}, a tree language has finite-algebra
recognition exactly when it is MSO definable.
For a supplied finite algebra and designated set, membership and emptiness
are effectively decidable. The statement uses finite input trees, not
arbitrary subsets of the full infinite successor tree.
\end{hbcorollary}

The MSO-to-automaton construction already aligns free-variable tracks,
enforces singleton flags for individuals, complements relative to legal
markings, and projects quantified tracks. Determinization then yields
the finite algebra. In the reverse direction, one set per algebra element
records the unique evaluation state at each node; the existing partition,
transition, and final-root clauses describe its accepting run.
No rank-changing erasure is introduced by this bridge.

WSkS's separate finite-set encoding over the infinite successor structure
requires the already proved canonical support/pruning construction of
\cref{sec:p09-wsks-canonical,sec:p09-wsks-pruning,sec:p09-wsks-logical-induction}.
This algebra treatment does not replace it with arbitrary bit erasure.
Full SkS requires its different infinite-tree machinery.

\begin{proofboundary}{def:p09-tw1968-unheld}{The original Thatcher--Wright presentation is still distinct}
The 1968 source named in N-c-015 remains a separate complete
pp.~57--81 obligation \parencite{ThatcherWright1968}.
Its native algebra conventions, closure and generalized Kleene
presentation, exact weak second-order language/coding, examples, and
proof architecture cannot be identified from an abstract with this
2008 treatment. The original journal body, its earlier RC1713 precursor,
and this modern algebra/MSO seam must retain their separate source
standing. No original-source admission follows from these new proofs.
\end{proofboundary}
